%------------------------------------------------------------------

Sources are organized into ten thematic areas. The per-area counts shown
in parentheses sum to 100, more than the 89 systematically-searched sources,
for two reasons: a source that substantively informs more than one area is
discussed under each and counted in both (here only LLMParser~\cite{llmparser}), and the
supplementary companion and unnumbered entries described below each carry
their own heading. Every systematically-searched source nonetheless appears
exactly once in the single-count tier classification of
the source-classification table of the thesis. Each
entry reports the acquisition channel, a summary, pro/con arguments, and
the usability verdict. Two kinds of entry appear in these appraisals
without counting toward the 89/93 review total: five companion entries,
close relatives discussed alongside their parent
source, and five further foundational or methodological citations
(Cover--Thomas, Gini/Yitzhaki, Lee et al.\ 2022, Inam et al.\ 2023, and
NIST SP~800-92), assessed for transparency but lying outside the
systematic search.

%---------------------------
\subsection{Benchmark Critiques and Log-Dataset Fitness (11 sources)}
\label{sec:bg-landscape}
%---------------------------

The benchmark-critique strand of the literature motivates the core 
question of this work. Loghub~\cite{loghub2023} consolidates 16 public log
datasets including HDFS, Hadoop, OpenStack, BGL, Thunderbird, and
others; several are repackaged HPC fault logs~\cite{oliner2007,%
schroeder2006} originally collected for reliability research, not
security evaluation. The AIT log datasets~\cite{ait2021have,%
ait2023maintainable} are produced by a model-driven simulation testbed
with configurable normal and attack behavior; three testbed
scenarios (Russell Mitchell, Santos and Wilson) are profiled in the thesis.
Provenance and endpoint datasets (ATLASv2~\cite{atlasv2}, the DARPA
TC streams CADETS and THEIA~\cite{darpa-tc-e3}) and the network-telemetry dataset
UWF-ZeekData22~\cite{uwfzeek22} are also included.

\slrpaper{Landauer, Skopik, Wurzenberger (2024) -- A Critical Review of Common Log Data Sets}
{Max Landauer, Florian Skopik, Markus Wurzenberger}
{Proc.\ ACM Softw.\ Eng.\ (FSE) 1:1354--1375, 2024}
{Empirically dissects HDFS, BGL, Thunderbird, OpenStack and Hadoop and
shows widespread redundancy, label leakage and non-sequential
anomalies. Operationalises redundancy via entropy and
near-duplicate-sequence analysis.
\begin{itemize}[nosep]
  \item \acq{Channel~1 (seed).}
  \item \datasets{HDFS, BGL, Thunderbird, OpenStack, Hadoop (the five
  Loghub benchmarks dissected; all in-corpus).}
  \item \pro{Highly relevant: first to systematically criticize
  benchmark reuse; directly precedents D1 metrics. Concrete numbers:
  OpenStack $>98.5\%$ overlap, Hadoop $>83\%$ duplicated sequences.
  The authors argue that high-capacity models may exploit benchmark
  simplicity and leakage rather than learn generalizable anomaly patterns,
  that reporting F1 alone is uninformative under severe class imbalance,
  and that the original data is often unavailable.}
  \item \con{The paper's recommendation of ADFA should be treated cautiously:
  ADFA was proposed over a decade earlier for a host-based IDS setting and
  does not directly address modern log-based anomaly detection.}
  \item \usmark{$\bigstar$ Primary: anchors D1, D3; provides canonical negative cases.}
\end{itemize}}

\slrpaper{Yu, Yao, Fu et al.\ (2024) -- Deep Learning or Classical Machine Learning?}
{Boxi Yu, Jiayi Yao, Qiuai Fu et al.}
{ICSE 2024}
{Shows KNN/DT/SLFN match or beat CNN/LogRobust/NeuralLog on five
benchmarks when tuned fairly; KNN is $\sim$1000$\times$ faster than
NeuralLog on Thunderbird.
\begin{itemize}[nosep]
  \item \acq{Channel~2 (forward snowball from \cite{landauer2024}).}
  \item \datasets{Five benchmarks: HDFS, BGL, Thunderbird (in-corpus),
  Spirit and Liberty (out-of-corpus).}
  \item \pro{Directly motivates the classical-ML probe family.
  On Thunderbird, KNN trains roughly three orders of magnitude faster than
  the deep baselines, illustrating that dataset simplicity and possible
  leakage can make detection unusually easy.}
  \item \con{Scope limited to five canonical datasets; no per-dataset
  decision rule.}
  \item \usmark{$\bigstar$ Primary: classical-ML baseline and saturation hypothesis.}
\end{itemize}}

\slrpaper{Le, Zhang (2022) -- Log-Based Anomaly Detection with Deep Learning: How Far Are We?}
{Van-Hoang Le, Hongyu Zhang}
{ICSE 2022, pp.\ 1356--1367}
{Demonstrates 1\% mislabeled logs can drop F1 by $>$35\% on
LogRobust/CNN. Quantifies parsing-error impact ($\sim$29\% F1
improvement for SVM when errors are fixed).
\begin{itemize}[nosep]
  \item \acq{Channel~2 (backward snowball from \cite{landauer2024}).}
  \item \datasets{HDFS, BGL, Thunderbird (in-corpus) and Spirit
  (out-of-corpus).}
  \item \pro{Provides the label-noise protocol (1\%, 5\%, 10\%)
  adopted in D3 sensitivity checks.}
  \item \con{Experiments on HDFS/BGL/Thunderbird/Spirit, the same
  legacy benchmarks the thesis questions.}
  \item \usmark{$\bigstar$ Primary: D3 label quality, NeuralLog probe family.}
\end{itemize}}

\slrpaper{Ali, Boufaied, Bianculli, Branco, Briand (2025) -- A Comprehensive Study of ML for Log-Based AD}
{Shan Ali, Chaima Boufaied, Domenico Bianculli, Paula Branco, Lionel Briand}
{Empir.\ Softw.\ Eng.\ 30(5):129, 2025}
{Evaluates supervised/semi-supervised ML along four axes: detection
accuracy, training time, prediction time, HP sensitivity. Argues
semi-supervised methods are underrepresented.
\begin{itemize}[nosep]
  \item \acq{Channel~2 (forward snowball from \cite{landauer2024}).}
  \item \datasets{HDFS, BGL, Thunderbird, Hadoop, OpenStack (in-corpus)
  and Spirit (out-of-corpus); Loghub benchmarks throughout.}
  \item \pro{Provides the multi-axis reporting template (Table~4).}
  \item \con{Still F1-centric; does not link results to dataset
  properties.}
  \item \usmark{$\bigstar$ Primary: defines how probe results are reported.}
\end{itemize}}

\slrpaper{Kenyon, Deka, Elizondo (2020) -- Are Public Intrusion Datasets Fit for Purpose?}
{A.\ Kenyon, L.\ Deka, D.\ Elizondo}
{Comput.\ Secur.\ 99:102022, 2020}
{Surveys intrusion-dataset landscape; identifies representativeness,
completeness, and labeling as recurring weaknesses.
\begin{itemize}[nosep]
  \item \acq{Channel~3 (SQ5).}
  \item \pro{Frames the fit-for-purpose question formalized by the
  thesis. Relative to earlier checklists it compares more datasets and places
  greater emphasis on provenance and dataset maintenance.}
  \item \con{Its scoring system is expert-opinion based rather than operational
  or computable, which is precisely the gap this thesis addresses.}
  \item \usmark{$\star$ Secondary: motivation and positioning.}
\end{itemize}}

\slrpaper{Garc\'ia G\'omez, Landauer et al.\ (2026) -- Collaborative Anomaly Detection: Comparative Analysis}
{Andr\'e Garc\'ia G\'omez, Max Landauer et al.}
{Future Generation Computer Systems 175:108090, 2025/2026}
{Extends the fit-for-purpose critique to collaborative log-AD; argues
HDFS/BGL were not designed for intrusion detection. Positions AIT~v2
as a more appropriate modern successor.
\begin{itemize}[nosep]
  \item \acq{Channel~1 (seed).}
  \item \datasets{HDFS and AIT-LDS (both in-corpus).}
  \item \pro{Directly supports AIT-LDS/AIT-ADS as calibration
  instrument and legacy benchmarks as negative cases.}
  \item \con{Evaluation remains expert-opinion-based scoring; does not
  produce a computable decision rule.}
  \item \usmark{$\bigstar$ Primary: justifies dataset roles.}
\end{itemize}}

\slrpaper{Gharib, Sharafaldin, Lashkari, Ghorbani (2016) -- An Evaluation Framework for Intrusion Detection Datasets}
{Amirhossein Gharib, Iman Sharafaldin, Ali Habibi Lashkari, Ali A.\ Ghorbani}
{ICISS 2016}
{Proposes 11 criteria for IDS datasets and scores several public
datasets.
\begin{itemize}[nosep]
  \item \acq{Channel~3 (SQ5).}
  \item \pro{A historical checklist-style framework whose criteria identify
  several relevant concerns; the core ideas remain useful.}
  \item \con{Largely a critique without an operational solution (binary criteria
  plus a diversity score), and now dated.}
  \item \usmark{$\star$ Secondary: historical framing.}
\end{itemize}}

\slrpaper{Shiravi, Shiravi, Tavallaee, Ghorbani (2012) -- Toward Developing a Systematic Approach to Generate Benchmark Datasets for IDS}
{Ali Shiravi, Hadi Shiravi, Mahbod Tavallaee, Ali A.\ Ghorbani}
{Comput.\ Secur.\ 31(3):357--374, 2012}
{Profile-based IDS dataset generation; introduces the ISCX dataset.
\begin{itemize}[nosep]
  \item \acq{Channel~3 (SQ5).}
  \item \pro{Early precedent for synthetic/generator-based benchmarks.}
  \item \con{Network-traffic focused; dataset now dated.}
  \item \usmark{$\circ$ Contextual: synthetic-benchmark lineage.}
\end{itemize}}

\slrpaper{Creech, Hu (2013) -- Generation of a New IDS Test Dataset: Time to Retire the KDD Collection}
{Gideon Creech, Jiankun Hu}
{IEEE WCNC 2013, pp.\ 4487--4492}
{Introduces ADFA dataset and argues KDD is no longer representative.
\begin{itemize}[nosep]
  \item \acq{Channel~3 (SQ5).}
  \item \pro{Seminal dataset-retirement argument; cited in \cite{landauer2024} reading
  notes.}
  \item \con{ADFA itself is now aged; only marginally applicable to
  modern log-based security analytics.}
  \item \usmark{$\circ$ Contextual: historical reference.}
\end{itemize}}

\slrpaper{Albert (2025) -- System Logs Anomaly Detection: Are We on the Right Path?}
{Ramona-Georgiana Albert}
{Survey/opinion, 2025}
{Position piece questioning progress in log-AD research.
\begin{itemize}[nosep]
  \item \acq{Channel~3 (SQ1).}
  \item \pro{Reinforces thesis motivation.}
  \item \con{Non-empirical.}
  \item \usmark{$\circ$ Contextual.}
\end{itemize}}

\slrpaper{He, Zhu, He, Lyu (2016) -- Experience Report: System Log Analysis for Anomaly Detection}
{Shilin He, Jieming Zhu, Pinjia He, Michael R.\ Lyu}
{ISSRE 2016, pp.\ 207--218}
{Early systematic comparison of classical AD methods; introduces much
of the now-standard HDFS session-based framing.
\begin{itemize}[nosep]
  \item \acq{Channel~2 (backward snowball from \cite{landauer2024}).}
  \item \datasets{HDFS, BGL (both in-corpus).}
  \item \pro{Foundational citation for the session-based HDFS framing.}
  \item \con{Helped entrench the two benchmarks whose reuse the thesis
  critiques.}
  \item \usmark{$\bigstar$ Primary context: motivation for the
  benchmark-critique strand.}
\end{itemize}}

%---------------------------
\subsection{Log Parsing and Its Downstream Effects (6 sources)}
\label{sec:bg-parsing}
%---------------------------

Most log-analytics methods do not operate directly on raw lines. A
preprocessing step called \emph{log parsing} converts each line into a
\emph{template} (a constant skeleton, e.g.~\code{Connection from <*>
closed by <*>}) and a parameter list, so that subsequent analyses can
operate on a finite event vocabulary. Drain~\cite{drain} is the most
widely used log parser in academic evaluations; its Drain3 extension
adds persistence and configurable masking but leaves the underlying
tree algorithm unchanged. Fu et~al.~\cite{fu2023parser} and Jiang et
al.~\cite{jiang2024loghub2} show that parser choice and template count
affect downstream AD performance independently of grouping accuracy,
a finding that motivates the parser-sensitivity robustness axis in the
thesis. Wu et~al.~\cite{wu2023logrep} compare classical message-count
vectors, template-ID sequences, and frozen contextual embeddings as
input representations for downstream AD.

\slrpaper{He, Zhu, Zheng, Lyu (2017) -- Drain: An Online Log Parsing Approach with Fixed Depth Tree}
{Pinjia He, Jieming Zhu, Zibin Zheng, Michael R.\ Lyu}
{ICWS 2017, pp.\ 33--40}
{Fixed-depth parse-tree streaming parser; best-in-class on
BGL/HDFS/Zookeeper (0.99 each); 51--81\% faster than prior online
parsers.
\begin{itemize}[nosep]
  \item \acq{Channel~1 (foundational seed).}
  \item \datasets{Five datasets: HDFS, BGL, HPC, Zookeeper, Proxifier
  (parsing-accuracy evaluation; HDFS/BGL in-corpus).}
  \item \pro{The de-facto parser in the field; detailed parsing-accuracy and
  runtime measurements are reported.}
  \item \con{Heuristic; accuracy drops on parameter-intensive
  templates (see \cite{jiang2024loghub2}).}
  \item \usmark{$\bigstar$ Primary: baseline parser for all D1--D5 measurements.}
\end{itemize}}

\slrpaper{Jiang, Liu, Huang et al.\ (2024) -- A Large-Scale Evaluation for Log Parsing Techniques: How Far Are We?}
{Zhihan Jiang, Jinyang Liu, Junjie Huang et al.}
{ISSTA 2024}
{Re-evaluates 15 parsers on Loghub-2.0 ($\sim$1000$\times$ larger
than Loghub-2k). All parsers drop in accuracy; Drain strongest for
grouping; semantic parsers stronger on token/template ID. 9/15 fail
12-hour timeout.
\begin{itemize}[nosep]
  \item \acq{Channel~3 (SQ2).}
  \item \datasets{Loghub-2.0 (14 large-scale datasets); counted in the
  census for HDFS, BGL, Thunderbird, Hadoop, OpenStack, HPC, Proxifier
  and Zookeeper.}
  \item \pro{Template-level metrics (FGA/FTA) and parameter-count
  sensitivity reused in D1. It provides a comprehensive account of Loghub-2.0 limitations and parser
  behavior.}
  \item \con{Focus is parsers, not downstream suitability; the
  coupling is the thesis's contribution.}
  \item \usmark{$\bigstar$ Primary: D1 metric basis, parser-sensitivity protocol.}
\end{itemize}}

\slrpaper{Zhang, Jiang, Zhang (2023) -- An Empirical Study of the Impact of Log Parsers on AD}
{Ying Fu, Meng Yan et al.}
{Empir.\ Softw.\ Eng.\ 28(6), 2023}
{Shows high parsing accuracy $\neq$ high AD performance. IPLoM and
Drain dominate downstream. Number of extracted templates matters
independently of accuracy.
\begin{itemize}[nosep]
  \item \acq{Channel~3 (SQ2).}
  \item \datasets{Six Loghub datasets (incl.\ HDFS, BGL, Thunderbird).}
  \item \pro{Evidence that preprocessing silently dominates reported
  scores. The study observes that parsing errors are occasionally beneficial and
  that heuristic parsers are the most suitable overall.}
  \item \con{Only 6 datasets; help-not-hurt mechanisms not fully
  unpacked.}
  \item \usmark{$\bigstar$ Primary: robustness-axis motivation.}
\end{itemize}}

\slrpaper{Khan, Shin, Bianculli, Briand (2024) -- Impact of Log Parsing on Deep Learning-Based Anomaly Detection}
{Zanis Ali Khan, Donghwan Shin, Domenico Bianculli, Lionel Briand}
{Empir.\ Softw.\ Eng., 2024}
{Confirms parser choice can flip relative DL model ranking.
\begin{itemize}[nosep]
  \item \acq{Channel~3 (SQ2).}
  \item \datasets{HDFS, Hadoop, BGL, OpenStack (in-corpus).}
  \item \pro{Reinforces D1/D2 parser-perturbation protocol.}
  \item \con{Narrow DL focus; overlaps with \cite{fu2023parser}.}
  \item \usmark{$\star$ Secondary.}
\end{itemize}}

\slrpaper{Le, Zhang (2021) -- Log-Based Anomaly Detection Without Log Parsing (NeuralLog)}
{Van-Hoang Le, Hongyu Zhang}
{ASE 2021}
{Bypasses the parser: raw logs $\to$ BERT $\to$ Transformer. On HDFS,
BGL, Thunderbird and Spirit F1 $\geq 0.96$.
\begin{itemize}[nosep]
  \item \acq{Channel~2 (second-order snowball from \cite{lezhang2022} via \cite{landauer2024}).}
  \item \datasets{HDFS, BGL, Thunderbird (in-corpus) and Spirit
  (out-of-corpus).}
  \item \pro{Shows semantic embedding carries signal beyond templates,
  relevant to D5. Correcting parsing errors improved SVM F1 by roughly 29\%.}
  \item \con{Parser-free pipelines still evaluated on the same legacy
  benchmarks.}
  \item \usmark{$\bigstar$ Primary: representation-sensitivity checks for D1 and D5.}
\end{itemize}}

\slrpaper{Ma, Chen, Kim, Chen, Wang (2024) -- LLMParser}
{Zeyang Ma, An Ran Chen, Dong Jae Kim et al.}
{ICSE 2024}
{Explores LLMs as log parsers; accuracy/cost trade-offs vs.\ Drain.
\begin{itemize}[nosep]
  \item \acq{Channel~3 (SQ2, SQ3).}
  \item \datasets{Full Loghub-2k suite (incl.\ HDFS, BGL, Thunderbird,
  OpenStack, Hadoop, Zookeeper, Proxifier; parsing benchmark, no
  end-to-end AD evaluation).}
  \item \pro{Enables optional LLM-parser class in parser-sensitivity
  check.}
  \item \con{Computational cost and prompt sensitivity may confound.}
  \item \usmark{$\star$ Secondary: optional parser class.}
\end{itemize}}

%---------------------------
\subsection{Anomaly Detection Methods (32 sources)}
\label{sec:bg-ad}
%---------------------------

Sequence-based deep learning (DeepLog~\cite{deeplog},
LogAnomaly~\cite{loganomaly}, LogRobust~\cite{logrobust},
LogFiT~\cite{logfit}) is the dominant baseline class in published
evaluations. Yu et~al.~\cite{yu2024dl} compare deep-learning sequence
detectors with classical machine-learning methods and report that
classical methods often match or exceed deep models, a finding which,
together with~\cite{landauer2024}, suggests deep models may be
exploiting dataset redundancy. Ali et~al.~\cite{ali2025} provide a
comprehensive multi-axis study under uniform evaluation conditions.
The earliest sequence-based detectors used $n$-gram and Markov models
over system-call sequences~\cite{forrest1996,hofmeyr1998}; the thesis
re-evaluates Markov-order estimation in the context of log templates as
the D2 probe. Friedberg et~al.~\cite{friedberg2015} address combined
event correlation for APT detection, foreshadowing D4; Garc\'ia
G\'omez et~al.~\cite{garcia2025} propose a collaborative AD framework
evaluated on AIT-LDS that informs parts of the D4 operationalization.

\subsubsection*{Foundational and Classical Methods}

\slrpaper{Forrest, Hofmeyr, Somayaji, Longstaff (1996) -- A Sense of Self for Unix Processes}
{S.\ Forrest, S.A.\ Hofmeyr, A.\ Somayaji, T.A.\ Longstaff}
{IEEE S\&P 1996, pp.\ 120--128}
{Seminal system-call sequence IDS via short-sequence windows.
\begin{itemize}[nosep]
  \item \acq{Channel~2 (backward snowball from \cite{landauer2024}).}
  \item \pro{Originator of the $n$-gram/Markov baseline the
  statistical/frequency probe uses.}
  \item \con{Pre-template era; cited for lineage.}
  \item \usmark{$\bigstar$ Primary: statistical probe lineage.}
\end{itemize}}

\slrpaper{Hofmeyr, Forrest, Somayaji (1998) -- Intrusion Detection Using Sequences of System Calls}
{S.A.\ Hofmeyr, S.\ Forrest, A.\ Somayaji}
{J.\ Comput.\ Secur.\ 6(3):151--180, 1998}
{Formalized distance measures over system-call sequences.
\begin{itemize}[nosep]
  \item \acq{Channel~2 (backward snowball from \cite{forrest1996}).}
  \item \pro{Historical anchor for Markov/$n$-gram baselines.}
  \item \con{Pre-template era.}
  \item \usmark{$\star$ Secondary.}
\end{itemize}}

\subsubsection*{Deep Sequence-Based AD (Probe Family Methods)}

\slrpaper{Du, Li, Zheng, Srikumar (2017) -- DeepLog}
{Min Du, Feifei Li, Guineng Zheng, Vivek Srikumar}
{CCS 2017, pp.\ 1285--1298}
{LSTM-based log-key + parameter-value sequence AD; HDFS F1
$\sim$96\%, trained on $<$1\% of normal sessions.
\begin{itemize}[nosep]
  \item \acq{Channel~1 (foundational seed).}
  \item \datasets{HDFS, OpenStack (both in-corpus).}
  \item \pro{The canonical deep sequential AD baseline; detailed per-dataset
  performance and online-update results are reported.}
  \item \con{\cite{landauer2024} argues that much of the reported gap is attributable to
  dataset leakage.}
  \item \usmark{$\bigstar$ Primary: sequence probe.}
\end{itemize}}

\slrpaper{Meng, Liu et al.\ (2019) -- LogAnomaly}
{Weibin Meng, Ying Liu et al.}
{IJCAI 2019, pp.\ 4739--4745}
{Sequential + quantitative AD with template2vec and attention-LSTM.
\begin{itemize}[nosep]
  \item \acq{Channel~2 (backward snowball from \cite{landauer2024}).}
  \item \datasets{HDFS, BGL (both in-corpus).}
  \item \pro{Representative semantic-sequence family member.}
  \item \con{Shares benchmark dependency with DeepLog.}
  \item \usmark{$\bigstar$ Primary: sequence probe.}
\end{itemize}}

\slrpaper{Zhang et al.\ (2019) -- LogRobust: Robust Log-Based Anomaly Detection on Unstable Log Data}
{Xu Zhang, Yong Xu et al.}
{ESEC/FSE 2019, pp.\ 807--817}
{Bi-LSTM with attention and semantic embeddings for log evolution.
\begin{itemize}[nosep]
  \item \acq{Channel~2 (backward snowball from \cite{landauer2024}).}
  \item \datasets{HDFS (in-corpus), plus a synthetic dataset derived from
  an industrial system; the original LogRobust evaluation does not use
  BGL.}
  \item \pro{Third canonical sequence baseline; central to
  leakage/label-noise analyses in \cite{lezhang2022}.}
  \item \con{The ``unstable log'' evaluation is partly circular if the
  underlying HDFS data already leaks.}
  \item \usmark{$\bigstar$ Primary: sequence probe.}
\end{itemize}}

\subsubsection*{Surveys of Log-AD Methods}

\slrpaper{Landauer, Onder et al.\ (2023) -- Deep Learning for Anomaly Detection in Log Data: A Survey}
{Max Landauer, Sebastian Onder et al.}
{2023}
{62-paper survey of DL for log-AD: architectures, representations,
evaluation practices.
\begin{itemize}[nosep]
  \item \acq{Channel~1 (seed).}
  \item \pro{A comprehensive DL landscape; the survey documents a heavy dependence
  on HDFS, BGL, Thunderbird and OpenStack and notes that public datasets
  are frequently reused.}
  \item \con{Survey-level; no new empirical data.}
  \item \usmark{$\bigstar$ Primary: landscape map for probe-family selection.}
\end{itemize}}

\slrpaper{Yadav, Kumar et al.\ (2020) -- A Survey on Log AD Using Deep Learning}
{Rakesh Bahadur Yadav, P Santosh Kumar et al.}
{2020}
{Earlier DL-for-LogAD survey.
\begin{itemize}[nosep]
  \item \acq{Channel~3 (SQ1).}
  \item \pro{Landscape coverage.}
  \item \con{Superseded by \cite{logdl_survey} and \cite{ali2025}.}
  \item \usmark{$\circ$ Contextual.}
\end{itemize}}

\subsubsection*{Additional AD Methods and Variants}

\slrpaper{Chen, Singh, Rahmani et al.\ (2025) -- $K^4$: Unsupervised Typicality Learning}
{Weicong Chen, Vikash Singh et al.}
{IEEE International Conference on High Performance Computing, Data, and Analytics (HiPC) 2025, pp.\ 96--107}
{Parser-free, unsupervised online LogAD.
\begin{itemize}[nosep]
  \item \acq{Channel~2d (forward snowball from Loghub, \cite{loghub2023}).}
  \item \datasets{Loghub benchmarks (HDFS, BGL, Thunderbird; in-corpus).}
  \item \pro{Supports unsupervised probe class; aligns with D3
  concerns.}
  \item \con{Classical benchmark evaluation.}
  \item \usmark{$\star$ Secondary.}
\end{itemize}}

\slrpaper{Chu, Wang et al.\ (2025) -- AD on Interleaved Log Data via Log-Entity Graph Mining}
{Guojun Chu, Jingyu Wang et al.}
{2025}
{Entity-graph perspective on interleaved logs.
\begin{itemize}[nosep]
  \item \acq{Channel~3 (SQ1).}
  \item \datasets{HDFS, BGL, Thunderbird (Loghub; in-corpus).}
  \item \pro{Directly related to D4 cross-component structure.}
  \item \con{Limited dataset scope.}
  \item \usmark{$\bigstar$ Primary: D4 related-work anchor.}
\end{itemize}}

\slrpaper{Fu, Lou et al.\ (2009) -- Execution Anomaly Detection in Distributed Systems}
{Qiang Fu, Jian-Guang Lou et al.}
{2009}
{Early pipeline for log-sequence AD on distributed systems.
\begin{itemize}[nosep]
  \item \acq{Channel~2 (backward snowball from \cite{logdl_survey}).}
  \item \datasets{Own pre-Loghub Hadoop-ecosystem collection (plus SILK);
  counted under the census's Hadoop row (the dataset-usage table of the thesis)
  although the collection predates the Loghub YARN release -- the same
  one-row-per-genre convention \cite{xu2009console}'s entry applies to the HDFS-origin
  logs.}
  \item \pro{Historical anchor for sequence-AD on real system logs.}
  \item \con{Pre-DL.}
  \item \usmark{$\circ$ Contextual.}
\end{itemize}}

\slrpaper{Chen, Liu et al.\ (2022) -- Experience Report: DL-Based System Log Analysis for AD}
{Zhuangbin Chen, Jinyang Liu et al.}
{2022}
{Comparative empirical report of DL-based log-AD methods.
\begin{itemize}[nosep]
  \item \acq{Channel~2 (backward snowball from \cite{logdl_survey}).}
  \item \datasets{HDFS, BGL (Loghub; in-corpus).}
  \item \pro{Reproducibility perspective on LogRobust/DeepLog.}
  \item \con{Summarises results on standard benchmarks.}
  \item \usmark{$\circ$ Contextual.}
\end{itemize}}

\slrpaper{FastLogAD -- Lin, Deng et al.\ (2024)}
{Yifei Lin, Hanqiu Deng et al.}
{IEEE AITest 2024 invited}
{Mask-guided pseudo-anomaly generation for discriminative AD.
\begin{itemize}[nosep]
  \item \acq{Channel~2d (forward snowball from Loghub, \cite{loghub2023}).}
  \item \datasets{HDFS, BGL, Thunderbird (Loghub; in-corpus).}
  \item \pro{Synthetic-anomaly idea relevant to calibration.}
  \item \con{Short paper; limited scope.}
  \item \usmark{$\circ$ Contextual.}
\end{itemize}}

\slrpaper{BERT-LogAnom -- Lu, An et al.\ (2026)}
{Xi Lu, Shufan An et al.}
{2026}
{BiLSTM + gated residual + dynamic thresholding on BERT embeddings.
\begin{itemize}[nosep]
  \item \acq{Channel~2d (forward snowball from Loghub, \cite{loghub2023}).}
  \item \datasets{HDFS, BGL, Thunderbird (Loghub; in-corpus).}
  \item \pro{Recent representation-driven AD baseline.}
  \item \con{Incremental over LogRobust/LogAnomaly lineage.}
  \item \usmark{$\circ$ Contextual.}
\end{itemize}}

\slrpaper{LogFormer -- Guo, Yang et al.\ (2024)}
{Hongcheng Guo, Jian Yang et al.}
{AAAI Conference on Artificial Intelligence 38(1), pp.\ 135--143, 2024}
{Pre-train + tune for log AD.
\begin{itemize}[nosep]
  \item \acq{Channel~2d (forward snowball from Loghub, \cite{loghub2023}).}
  \item \datasets{BGL, Thunderbird (in-corpus) and Spirit, Liberty
  (out-of-corpus).}
  \item \pro{Pre-train/tune pipeline relevant to LM probes.}
  \item \con{Self-reported overlap reduces independent contribution.}
  \item \usmark{$\circ$ Contextual.}
\end{itemize}}

\slrpaper{LayerLog -- Zhang, Wang et al.\ (2023)}
{Chunkai Zhang, Xinyu Wang et al.}
{2023}
{Hierarchical semantic AD on log sequences.
\begin{itemize}[nosep]
  \item \acq{Channel~3 (SQ1).}
  \item \datasets{HDFS, BGL (Loghub; in-corpus).}
  \item \pro{Recent hierarchical variant for sequence family.}
  \item \con{Little new validation on D4/D5.}
  \item \usmark{$\star$ Secondary.}
\end{itemize}}

\slrpaper{Ott, Bogatinovski et al.\ (2021) -- Robust and Transferable AD with Pre-Trained LMs}
{Harold Ott, Jasmin Bogatinovski et al.}
{2021}
{Pre-trained-LM-based transferable AD across datasets.
\begin{itemize}[nosep]
  \item \acq{Channel~3 (SQ1).}
  \item \pro{Motivates D5 representation-sensitivity checks.}
  \item \con{Transferability claims need re-verification.}
  \item \usmark{$\star$ Secondary.}
\end{itemize}}

\slrpaper{Jia, Cai et al.\ (2023) -- Robust and Transferable Log-Based AD}
{Peng Jia, Shaofeng Cai et al.}
{2023}
{Variant of transferable LM-based AD.
\begin{itemize}[nosep]
  \item \acq{Channel~3 (SQ1).}
  \item \pro{Additional transferability baseline.}
  \item \con{Similar scope; standard datasets.}
  \item \usmark{$\star$ Secondary.}
\end{itemize}}

\slrpaper{Yang, Chen et al.\ (2021) -- Semi-Supervised Log AD via Probabilistic Label Estimation}
{Lin Yang, Junjie Chen et al.}
{2021}
{Probabilistic label estimation for partially labeled logs.
\begin{itemize}[nosep]
  \item \acq{Channel~3 (SQ1).}
  \item \datasets{HDFS, BGL (Loghub; in-corpus).}
  \item \pro{Relevant to D3 label quality and noise tolerance.}
  \item \con{Specific estimation assumptions.}
  \item \usmark{$\star$ Secondary.}
\end{itemize}}

\slrpaper{Wurzenberger, Skopik et al.\ (2017) -- Incremental Clustering for Semi-Supervised AD on Log Data}
{Markus Wurzenberger, Florian Skopik et al.}
{2017}
{Incremental clustering for semi-supervised log AD.
\begin{itemize}[nosep]
  \item \acq{Channel~1 (seed).}
  \item \datasets{AIT / AECID log data (in-corpus AIT-LDS family).}
  \item \pro{Semi-supervised probe candidate.}
  \item \con{Specific to AECID family.}
  \item \usmark{$\star$ Secondary.}
\end{itemize}}

\slrpaper{Beck, Landauer et al.\ (2024) -- Semi-Supervised Configuration and Optimization of AD}
{Viktor Beck, Max Landauer et al.}
{2024}
{Automated tuning for semi-supervised log-AD.
\begin{itemize}[nosep]
  \item \acq{Channel~1 (seed).}
  \item \datasets{AIT-LDS / AECID scenarios (incl.\ Russell Mitchell,
  Harrison, Fox; the AIT-LDS family is in-corpus).}
  \item \pro{Addresses HP-sensitivity axis from \cite{ali2025}.}
  \item \con{Ties to AECID stack.}
  \item \usmark{$\star$ Secondary.}
\end{itemize}}

\slrpaper{Alam, Kifayat et al.\ (2024) -- SXAD: Shapely eXplainable AI-Based AD Using Log Data (sic)}
{Kashif Alam, Kashif Kifayat et al.}
{2024}
{Shapley-value explainability on log-AD outputs.
\begin{itemize}[nosep]
  \item \acq{Channel~3 (SQ1).}
  \item \datasets{HDFS (Loghub; in-corpus).}
  \item \pro{Supporting citation for interpretation class.}
  \item \con{Post-hoc explainability, not dataset-centric.}
  \item \usmark{$\circ$ Contextual.}
\end{itemize}}

\slrpaper{Wang, Zhao et al.\ (2020) -- Root-Cause Metric Location for Microservices via Log AD}
{Lingzhi Wang, Nengwen Zhao et al.}
{2020}
{Microservice RCA via log-AD signals.
\begin{itemize}[nosep]
  \item \acq{Channel~3 (SQ1).}
  \item \pro{Supporting citation for RCA application class.}
  \item \con{Metrics component outside thesis scope.}
  \item \usmark{$\circ$ Contextual.}
\end{itemize}}

\slrpaper{Remaining contextual AD sources}
{Various authors}
{Various venues}
{Narrow-domain AD papers: access-log AD~\cite{accessloganom}, deep-log modeling~\cite{advlog_deeplogmodel}, application-log AD (\cite{applogdata}, MSc thesis), Hadoop-ecosystem AD~\cite{bigloghadoop}, server-log case study~\cite{serverlog_casestudy}, NLP-based log AD~\cite{nlplog}, federated
log AD~\cite{fedlightweight,fedlog_openchallenges}, physical-access-control AD~\cite{behavioraccess}, and
smart-contract AD~\cite{logsentry}.
\begin{itemize}[nosep]
  \item \acq{Channel~1~\cite{fedlog_openchallenges,behavioraccess}; Channel~2d (\cite{advlog_deeplogmodel,logsentry} -- forward
  snowball from Loghub, \cite{loghub2023}); Channel~3/SQ1 (remainder).}
  \item \datasets{The standard Loghub benchmarks, per source: \cite{advlog_deeplogmodel} (HDFS,
  BGL, Thunderbird, Spirit), \cite{applogdata} (HDFS, BGL), \cite{nlplog} (BGL), \cite{fedlog_openchallenges} (HDFS),
  \cite{fedlightweight} (HDFS, Thunderbird); a Hadoop-ecosystem dataset~\cite{bigloghadoop}; and
  domain-specific collections outside the ranked record datasets (web
  access logs \cite{accessloganom}, server logs \cite{serverlog_casestudy}, physical-access logs \cite{behavioraccess},
  smart-contract logs \cite{logsentry}). The Loghub-benchmark uses are tallied
  individually in the dataset-usage table of the thesis; the domain-specific
  collections are not counted there.}
  \item \pro{Illustrate breadth of log-AD application; \cite{fedlog_openchallenges,behavioraccess} from
  the seeds provide AECID cross-domain evidence.}
  \item \con{Narrow settings; \cite{logsentry} is outside scope.}
  \item \usmark{$\circ$ Contextual (\cite{fedlog_openchallenges,behavioraccess} = $\star$ Secondary).
  \cite{logsentry} = $\times$ Excluded.}
\end{itemize}}

%---------------------------
\subsection{LLM-Based Log Analysis (12 sources)}
%---------------------------

LLM-based approaches to log analysis have grown rapidly since 2023.
LogGPT~\cite{loggpt}, LogFiT~\cite{logfit}, and the survey by Akhtar
et~al.~\cite{akhtar2025} cover the landscape; Ma et~al.~\cite{llmparser}
document prompt sensitivity and hallucinations in LLM-based log parsing.
Wang et~al.~\cite{wang2025} provide a broader AI-driven log-analysis
survey.

\slrpaper{Qi, Huang, Luan et al.\ (2023) -- LogGPT: Exploring ChatGPT for Log-Based Anomaly Detection}
{Jiaxing Qi, Shaohan Huang, Zhongzhi Luan et al.}
{IEEE HPCC/DSS/SmartCity 2023}
{First ChatGPT-based log-AD framework (zero/few-shot) on BGL and
Spirit subset. Identifies prompt sensitivity, limited window, high
FP rate, hallucination.
\begin{itemize}[nosep]
  \item \acq{Channel~3 (SQ3).}
  \item \datasets{BGL (in-corpus) and a Spirit subset (out-of-corpus).}
  \item \pro{Enumerates failure modes the LLM-interpretation rubric
  (future work, the future-work section of the thesis) operationalises.
  Its main contributions are the prompt-design exploration and
  interpretability analysis; its key weaknesses are a high false-positive
  rate and hallucination.}
  \item \con{Sampled Spirit subset limits statistical reliability.}
  \item \usmark{$\bigstar$ Primary: anchors the LLM-interpretation prompt-variant rubric.}
\end{itemize}}

\slrpaper{Akhtar, Khan, Parkinson -- LLM-Based Event Log Analysis Techniques: A Survey}
{Siraaj Akhtar, Saad Khan, Simon Parkinson}
{arXiv:2502.00677 (GL)}
{Survey of LLM-based log-analysis approaches.
\begin{itemize}[nosep]
  \item \acq{Channel~3 (SQ3).}
  \item \pro{Landscape citation for the LLM-interpretation strand.}
  \item \con{Survey-level; does not measure semantic sufficiency.}
  \item \usmark{$\bigstar$ Primary: LLM-interpretation landscape.}
\end{itemize}}

\slrpaper{Wang, Gan, Yu (2025) -- AI-Driven Log Analysis: Advances and Challenges}
{Yongheng Wang, Wensheng Gan, Philip S.\ Yu}
{IEEE BigData 2025}
{Broader AI-driven log analysis survey.
\begin{itemize}[nosep]
  \item \acq{Channel~3 (SQ3).}
  \item \pro{Broader coverage including non-LLM AI.}
  \item \con{High-level; no new data.}
  \item \usmark{$\circ$ Contextual.}
\end{itemize}}

\slrpaper{Almod\'ovar, Sabrina, Karimi, Azad (2024) -- LogFiT}
{Crispin Almodovar, Fariza Sabrina et al.}
{IEEE Trans.\ Netw.\ Serv.\ Manag., 2024}
{Fine-tuned LM for log AD.
\begin{itemize}[nosep]
  \item \acq{Channel~2d (forward snowball from Loghub, \cite{loghub2023}).}
  \item \datasets{HDFS, BGL, Thunderbird (Loghub; all in-corpus).}
  \item \pro{Fine-tuned-LM branch of semantic probes.}
  \item \con{Fine-tuning on the legacy benchmarks inherits their
  documented issues.}
  \item \usmark{$\star$ Secondary.}
\end{itemize}}

\slrpaper{Yang, Harris (2025) -- LogLLaMA}
{Zhuoyi Yang, Ian G.\ Harris}
{arXiv:2503.14849 (GL)}
{LLaMA2 fine-tuned on normal logs + RL for AD on BGL/TB/HDFS.
\begin{itemize}[nosep]
  \item \acq{Channel~1 (foundational seed).}
  \item \datasets{BGL, Thunderbird, HDFS (Loghub; all in-corpus).}
  \item \pro{Modern LLM baseline for the interpretation strand;
  representative recent LLM-based detector.}
  \item \con{Preprint; canonical benchmarks only.}
  \item \usmark{$\circ$ Contextual.}
\end{itemize}}

\slrpaper{Guan, Cao et al.\ (2025) -- LogLLM}
{Wei Guan, Jian Cao et al.}
{arXiv preprint (GL)}
{BERT + LLaMA with projector; regex preprocessing instead of parser.
\begin{itemize}[nosep]
  \item \acq{Channel~2d (forward snowball from Loghub, \cite{loghub2023}).}
  \item \datasets{HDFS, BGL, Thunderbird (in-corpus) and Liberty
  (out-of-corpus).}
  \item \pro{Parser-free LLM pipeline relevant to D1/D5 sensitivity.}
  \item \con{Overlaps with LogFiT/LogLLaMA family.}
  \item \usmark{$\star$ Secondary.}
\end{itemize}}

\slrpaper{Huang, Zhang et al.\ -- LogRules}
{Xin Huang, Ting Zhang et al.}
{Preprint (GL)}
{Rule-augmented LLM for log analysis.
\begin{itemize}[nosep]
  \item \acq{Channel~2d (forward snowball from Loghub, \cite{loghub2023}).}
  \item \datasets{Full Loghub-2k suite: HDFS, BGL, Thunderbird,
  OpenStack, Hadoop (in-corpus) and Spirit (out-of-corpus).}
  \item \pro{Prompt-engineering variant useful for the LLM prompt-variant
  axis.}
  \item \con{Limited reproducibility data.}
  \item \usmark{$\circ$ Contextual.}
\end{itemize}}

\slrpaper{Reis, Areias et al.\ -- LLM Framework for Log Sequence AD}
{Jo\~ao Reis, Miguel Areias et al.}
{Progress in Artificial Intelligence (EPIA), LNCS, Springer, 2025, pp.\ 324--334}
{Generic LLM framework for log-sequence AD.
\begin{itemize}[nosep]
  \item \acq{Channel~2d (forward snowball from Loghub, \cite{loghub2023}).}
  \item \datasets{Loghub benchmarks (HDFS, BGL, Thunderbird; in-corpus).}
  \item \pro{Landscape coverage.}
  \item \con{Limited novelty.}
  \item \usmark{$\circ$ Contextual.}
\end{itemize}}

\slrpaper{Zhu, Cui et al.\ (2025) -- LLM with Semantic Compression for Log AD}
{Xiangyu Zhu, Wenhua Cui et al.}
{2025}
{Semantic compression to handle long-context log sequences for LLMs.
\begin{itemize}[nosep]
  \item \acq{Channel~2d (forward snowball from Loghub, \cite{loghub2023}).}
  \item \datasets{Loghub benchmarks (HDFS, BGL, Thunderbird; in-corpus).}
  \item \pro{Addresses LLM context-window limitation identified in
  \cite{loggpt}.}
  \item \con{Does not evaluate semantic sufficiency of underlying
  data.}
  \item \usmark{$\circ$ Contextual.}
\end{itemize}}

\slrpaper{Hadadi, Xu et al.\ (2025) -- LLM Meets ML: Data-Efficient AD on Unstable Logs}
{Fatemeh Hadadi, Qinghua Xu et al.}
{2025}
{Combines classical ML with LLM signals.
\begin{itemize}[nosep]
  \item \acq{Channel~3 (SQ1, SQ3).}
  \item \datasets{HDFS (unstable-log AD benchmark; in-corpus).}
  \item \pro{Relevant to D3 noise tolerance.}
  \item \con{Narrow scope.}
  \item \usmark{$\circ$ Contextual.}
\end{itemize}}

\slrpaper{Yu, Zhang et al.\ (2024) -- What Makes a High-Quality Training Dataset for LLMs}
{Xiao Yu, Zexian Zhang et al.}
{ASE 2024, pp.\ 656--668}
{Developer-survey study of LLM-training-data quality practices.
\begin{itemize}[nosep]
  \item \acq{Channel~3 (SQ4).}
  \item \pro{Practitioner views on DQ for LLMs; complements
  Kapoor/Narayanan for the LLM-interpretation framing.}
  \item \con{Not log-specific.}
  \item \usmark{$\star$ Secondary.}
\end{itemize}}

\slrpaper{Ma, Chen, Kim, Chen, Wang (2024) -- LLMParser}
{Zeyang Ma, An Ran Chen, Dong Jae Kim et al.}
{ICSE 2024}
{Explores LLMs as log parsers; accuracy/cost trade-offs vs.\ Drain.
\begin{itemize}[nosep]
  \item \acq{Channel~3 (SQ2, SQ3).}
  \item \datasets{Full Loghub-2k suite (incl.\ HDFS, BGL, Thunderbird,
  OpenStack, Hadoop, Zookeeper, Proxifier; parsing benchmark, no
  end-to-end AD evaluation).}
  \item \pro{Enables optional LLM-parser class; documents prompt
  sensitivity and hallucinations.}
  \item \con{Cost and prompt sensitivity may confound comparisons.}
  \item \usmark{$\star$ Secondary: optional parser class.}
\end{itemize}}

%---------------------------
\subsection{Data-Quality Frameworks (10 sources)}
\label{sec:bg-dq}
%---------------------------

The Wang--Strong taxonomy~\cite{wang1996} organises data-quality
dimensions into intrinsic (accuracy, consistency), contextual
(timeliness, completeness), representational (interpretability,
conciseness) and accessibility categories. Subsequent
surveys~\cite{batini2009,gong2023,zhou2024} consolidate the
literature. The METRIC framework~\cite{schwabe2024} specialises DQ
dimensions for medical AI datasets, providing a methodological
precedent for log-domain instantiation: dimensions are not invented
from scratch but selected from an established catalog and adapted to
domain-specific failure modes. This thesis follows the same approach.
Kapoor and Narayanan~\cite{kapoor2023} provide a taxonomy of eight
leakage types in ML-based science, used in this thesis to organise the
D3 and D4 robustness checks.

\slrpaper{Wang, Strong (1996) -- Beyond Accuracy: What Data Quality Means to Data Consumers}
{Richard Y.\ Wang, Diane M.\ Strong}
{J.\ Manag.\ Inf.\ Syst.\ 12(4):5--33, 1996}
{Classical DQ taxonomy: intrinsic, contextual, representational,
accessibility.
\begin{itemize}[nosep]
  \item \acq{Channel~3 (SQ4).}
  \item \pro{Root framework; every D1--D5 dimension is anchored here.}
  \item \con{General-purpose; log specialization is the thesis's job.}
  \item \usmark{$\bigstar$ Primary: methodological anchor.}
\end{itemize}}

\slrpaper{Schwabe, Becker, Seyferth, Kla\ss, Sch\"affter (2024) -- METRIC-Framework for Assessing DQ for Trustworthy AI in Medicine}
{Daniel Schwabe, Katinka Becker et al.}
{npj Digit.\ Med.\ 7(1):203, 2024}
{Domain-specialised DQ framework via PRISMA review for medical AI.
\begin{itemize}[nosep]
  \item \acq{Channel~3 (SQ4).}
  \item \pro{Structural template: shows how general DQ $\to$
  domain-specific dimensions via SLR.}
  \item \con{Medical-specific; domain details not transferable.}
  \item \usmark{$\bigstar$ Primary: methodological template.}
\end{itemize}}

\slrpaper{Zhou, Tu, Sha, Ding, Chen (2024) -- Survey on DQ Dimensions and Tools for ML}
{Yuhan Zhou, Fengjiao Tu et al.}
{arXiv:2406.19614; IEEE AITest 2024 invited (GL)}
{Catalogs DQ dimensions and tools for ML.
\begin{itemize}[nosep]
  \item \acq{Channel~3 (SQ4).}
  \item \pro{Expanded DQ vocabulary used in D1--D5 anchoring.}
  \item \con{No log-specific instantiation.}
  \item \usmark{$\bigstar$ Primary: DQ vocabulary.}
\end{itemize}}

\slrpaper{Gong, Liu, Xue, Li, Meng (2023) -- Survey on Dataset Quality in ML}
{Youdi Gong, Guangzhen Liu et al.}
{Inf.\ Softw.\ Technol.\ 162:107268, 2023}
{Lifecycle-oriented DQ survey.
\begin{itemize}[nosep]
  \item \acq{Channel~3 (SQ4).}
  \item \pro{Lifecycle angle on DQ.}
  \item \con{General-purpose.}
  \item \usmark{$\star$ Secondary.}
\end{itemize}}

\slrpaper{Kapoor, Narayanan (2023) -- Leakage and the Reproducibility Crisis in ML-Based Science}
{Sayash Kapoor, Arvind Narayanan}
{Patterns 4(9):100804, 2023}
{Formal 8-category leakage taxonomy for ML research.
\begin{itemize}[nosep]
  \item \acq{Channel~3 (SQ7).}
  \item \pro{Directly structures the leakage checks; each category
  maps to a log-data check.}
  \item \con{General ML; log instantiation contributed by thesis.}
  \item \usmark{$\bigstar$ Primary: leakage taxonomy backbone.}
\end{itemize}}

\slrpaper{Jayawardene, Sadiq et al.\ -- An Analysis of Data Quality Dimensions}
{Vimukthi Jayawardene, Shazia Sadiq et al.}
{Technical report}
{Dimension-level DQ analysis.
\begin{itemize}[nosep]
  \item \acq{Channel~3 (SQ4).}
  \item \pro{Supplementary reference.}
  \item \con{Not widely cited.}
  \item \usmark{$\circ$ Contextual.}
\end{itemize}}

\slrpaper{Batini, Cappiello, Francalanci, Maurino (2009) -- Methodologies for DQ Assessment and Improvement}
{Carlo Batini, Cinzia Cappiello et al.}
{ACM Comput.\ Surv.\ 41(3):1--52, 2009}
{Comprehensive DQ methodology survey.
\begin{itemize}[nosep]
  \item \acq{Channel~3 (SQ4).}
  \item \pro{Broad coverage; good for placing D1--D5 in a well-known
  landscape.}
  \item \con{Pre-ML era.}
  \item \usmark{$\bigstar$ Primary: DQ framing.}
\end{itemize}}

\slrpaper{Picard, Chapdelaine et al.\ (2020) -- Ensuring Dataset Quality for ML Certification}
{S.\ Picard, C.\ Chapdelaine et al.}
{2020}
{Certification-oriented DQ criteria.
\begin{itemize}[nosep]
  \item \acq{Channel~3 (SQ4).}
  \item \pro{Certification angle useful for reasoning about datasets
  whose quality cannot be conclusively judged either way.}
  \item \con{Embedded/safety-critical ML, not security analytics.}
  \item \usmark{$\circ$ Contextual.}
\end{itemize}}

\slrpaper{Chen, Chen et al.\ (2021) -- Data Evaluation and Enhancement for Quality Improvement of ML}
{Haihua Chen, Jiangping Chen et al.}
{2021}
{DQ improvement techniques.
\begin{itemize}[nosep]
  \item \acq{Channel~3 (SQ4).}
  \item \pro{Improvement-side perspective.}
  \item \con{Out of scope for assessment.}
  \item \usmark{$\circ$ Contextual.}
\end{itemize}}

\slrpaper{Zhao, Chen et al.\ (2020) -- On the Role of Dataset Quality and Heterogeneity in Model Confidence}
{Yuan Zhao, Jiasi Chen et al.}
{2020}
{Empirical study of DQ and heterogeneity effects on model confidence.
\begin{itemize}[nosep]
  \item \acq{Channel~3 (SQ4).}
  \item \pro{Supports the thesis's core claim that measurable dataset
  properties drive model behavior.}
  \item \con{Vision/tabular-oriented; log reinterpretation needed.}
  \item \usmark{$\star$ Secondary.}
\end{itemize}}

%---------------------------
\subsection{Statistical and Information-Theoretic Tools (9 sources)}
%---------------------------

\slrpaper{Shannon (1948) -- A Mathematical Theory of Communication}
{Claude E.\ Shannon}
{Bell Syst.\ Tech.\ J.\ 27(3):379--423, 1948}
{Defines entropy and mutual information.
\begin{itemize}[nosep]
  \item \acq{Channel~2 (backward snowball from \cite{landauer2024}).}
  \item \pro{Foundation for D1 entropy and D2 mutual-information
  metrics.}
  \item \usmark{$\bigstar$ Primary: foundational.}
\end{itemize}}

\slrpaper{Cover, Thomas (2006) -- Elements of Information Theory, 2nd ed.}
{Thomas M.\ Cover, Joy A.\ Thomas}
{Wiley-Interscience, 2006}
{Canonical textbook for MI, conditional MI, transition entropy.
\begin{itemize}[nosep]
  \item \acq{Channel~3 (SQ12 / supplementary).}
  \item \pro{Reference for D2 metrics.}
  \item \usmark{$\bigstar$ Primary: textbook reference.}
\end{itemize}}

\slrpaper{Gini (1912); Yitzhaki, Schechtman (2013) -- Gini Coefficient}
{C.\ Gini; S.\ Yitzhaki, E.\ Schechtman}
{Book / Springer 2013}
{Classical inequality measure.
\begin{itemize}[nosep]
  \item \acq{Channel~3 (SQ13 / supplementary).}
  \item \pro{Complements Shannon entropy in D1.}
  \item \usmark{$\bigstar$ Primary: D1 metric.}
\end{itemize}}

\slrpaper{Tangirala (2020) -- Evaluating the Impact of Gini Index and Information Gain on Classification}
{Suryakanthi Tangirala}
{2020}
{Empirical Gini vs.\ information-gain comparison.
\begin{itemize}[nosep]
  \item \acq{Channel~3 (SQ13).}
  \item \pro{Sanity check for D1 Gini choice.}
  \item \con{Decision-tree specific.}
  \item \usmark{$\circ$ Contextual.}
\end{itemize}}

\slrpaper{Papapetrou, Kugiumtzis (2013) -- Markov Chain Order Estimation with Conditional MI}
{Maria Papapetrou, Dimitris Kugiumtzis}
{Physica A 392(7):1593--1601, 2013}
{Conditional-MI-based Markov-order estimation.
\begin{itemize}[nosep]
  \item \acq{Channel~3 (SQ12).}
  \item \pro{Direct method reference for D2 Markov-order analysis.}
  \item \con{Small-scale; scalability on log datasets must be
  verified.}
  \item \usmark{$\bigstar$ Primary: D2 metric.}
\end{itemize}}

\slrpaper{Peng, Wang, Deng (2023) -- Near-Duplicate Sequence Search at Scale for LLM Memorization Evaluation}
{Zhencan Peng, Zhizhi Wang, Dong Deng}
{2023}
{Scalable MinHash/LSH near-duplicate search for LLM sets.
\begin{itemize}[nosep]
  \item \acq{Channel~3 (SQ11).}
  \item \pro{Scalable MinHash/LSH back-end for D1 near-duplicate
  detection.}
  \item \con{LLM-corpus setting; parameter tuning for log sequences
  needed.}
  \item \usmark{$\bigstar$ Primary: D1 implementation.}
\end{itemize}}

\slrpaper{Lee et al.\ (2022) -- Deduplicating Training Data Makes Language Models Better}
{Katherine Lee, Daphne Ippolito et al.}
{ACL 2022}
{Empirical impact of deduplication; popularises MinHash/LSH.
\begin{itemize}[nosep]
  \item \acq{Channel~3 (SQ11).}
  \item \pro{Precedent for D1 dedup methodology.}
  \item \con{NLP-specific.}
  \item \usmark{$\bigstar$ Primary: methodology precedent.}
\end{itemize}}

\slrpaper{Wurzenberger, H\"old, Landauer, Skopik (2024) -- Statistical Properties of Variables in Log Data}
{Markus Wurzenberger, Georg H\"old, Max Landauer, Florian Skopik}
{Comput.\ Secur.\ 137:103631, 2024}
{Statistical-property analysis of log-parameter values for AD.
\begin{itemize}[nosep]
  \item \acq{Channel~1 (seed).}
  \item \datasets{AIT-LDS / AECID testbed scenarios (in-corpus).}
  \item \pro{Direct basis for parameter-variability metrics in D1 and
  parameter-context coupling in D5.}
  \item \con{Single-family; cross-validation on non-AECID-parsed data
  needed.}
  \item \usmark{$\bigstar$ Primary: D1/D5 metrics.}
\end{itemize}}

\slrpaper{Kenyon, Deka et al.\ -- Characterising Payload Entropy in Packet Flows}
{Anthony Kenyon, Lipika Deka et al.}
{Journal}
{Entropy baselines for network payloads.
\begin{itemize}[nosep]
  \item \acq{Channel~3 (SQ13).}
  \item \pro{Supporting entropy-methodology reference.}
  \item \con{Packet focus, not log.}
  \item \usmark{$\circ$ Contextual.}
\end{itemize}}

%---------------------------
\subsection{Provenance, CTI, IDS Methodology and Cross-Component Structure (8 sources)}
%---------------------------

\slrpaper{Friedberg, Skopik, Settanni, Fiedler (2015) -- Combating APTs: Event Correlation to Incident Detection}
{Ivo Friedberg, Florian Skopik, Giuseppe Settanni, Roman Fiedler}
{Comput.\ Secur.\ 48:35--57, 2015}
{Cross-component event-correlation for APT detection.
\begin{itemize}[nosep]
  \item \acq{Channel~1 (seed).}
  \item \datasets{AIT / AECID event logs (in-corpus AIT-LDS family).}
  \item \pro{Direct precedent for D4 motif-stability; provides the
  motif-mining approach.}
  \item \con{APT-focused; generalization to RCA needed.}
  \item \usmark{$\bigstar$ Primary: D4 methodology.}
\end{itemize}}

\slrpaper{Agrawal, Imieli\'nski, Swami (1993) -- Mining Association Rules}
{R.\ Agrawal, T.\ Imieli\'nski, A.\ Swami}
{ACM SIGMOD 1993, pp.\ 207--216}
{Classical support/confidence for association-rule mining.
\begin{itemize}[nosep]
  \item \acq{Channel~2 (backward snowball from \cite{friedberg2015}).}
  \item \pro{Support/confidence measures used in D4 motif stability.}
  \item \con{Classical; log adaptation required.}
  \item \usmark{$\bigstar$ Primary: D4 metrics.}
\end{itemize}}

\slrpaper{Zipperle, Gottwalt et al.\ (2023) -- Provenance-Based IDS: A Survey}
{Michael Zipperle, Florian Gottwalt et al.}
{2023}
{Survey of provenance-graph IDS.
\begin{itemize}[nosep]
  \item \acq{Channel~3 (SQ6).}
  \item \pro{Justifies including the DARPA TC provenance datasets
  (CADETS, THEIA) and motivates D4 as a domain-specific dimension.}
  \item \con{Provenance-graph focus; log adaptation needed.}
  \item \usmark{$\bigstar$ Primary: D4 related work.}
\end{itemize}}

\slrpaper{Inam, Chen, Goyal et al.\ (2023) -- SoK: Provenance of System Intrusions}
{M.A.\ Inam, Y.\ Chen, A.\ Goyal et al.}
{IEEE S\&P 2023, pp.\ 2620--2638}
{Systematization of knowledge on audit-based provenance IDS.
\begin{itemize}[nosep]
  \item \acq{Channel~3 (SQ6 / citation follow-up from ATLASv2).}
  \item \pro{Sets state of the art in provenance IDS; supports D4.}
  \item \con{Provenance focus; log-AD is the fallback when D4 is
  uncertain.}
  \item \usmark{$\bigstar$ Primary: D4 SoK.}
\end{itemize}}

\slrpaper{Strom, Applebaum et al.\ (2018) -- MITRE ATT\&CK: Design and Philosophy}
{Blake E.\ Strom, Andy Applebaum et al.}
{MITRE Tech.\ Report (GL)}
{Design rationale of ATT\&CK.
\begin{itemize}[nosep]
  \item \acq{Channel~3 (SQ9).}
  \item \pro{Anchor for CTI/ATT\&CK application class.}
  \item \con{Not peer-reviewed; version pinning needed.}
  \item \usmark{$\bigstar$ Primary: CTI application class.}
\end{itemize}}

\slrpaper{Applebaum, Miller et al.\ (2016) -- Intelligent, Automated Red Team Emulation (CALDERA)}
{Andy Applebaum, Doug Miller et al.}
{2016}
{Automated adversary emulation.
\begin{itemize}[nosep]
  \item \acq{Channel~3 (SQ9).}
  \item \pro{Supports adversary-emulation justification.}
  \item \con{Tool paper; contextual.}
  \item \usmark{$\circ$ Contextual.}
\end{itemize}}

\slrpaper{Landauer, Skopik et al.\ (2025) -- A Review of Time-Series Analysis for Cyber Security Analytics}
{Max Landauer, Florian Skopik et al.}
{2025}
{Survey of time-series methods for cyber-security analytics.
\begin{itemize}[nosep]
  \item \acq{Channel~1 (seed).}
  \item \pro{Broadens D2 temporal methodology.}
  \item \con{Not log-specific throughout.}
  \item \usmark{$\star$ Secondary: D2 breadth reference.}
\end{itemize}}

\slrpaper{Hus\'ak, Bajto\v{s} et al.\ (2020) -- Predictive Cyber Situational Awareness via Sequential Rule Mining}
{Martin Hus\'ak, Tom\'a\v{s} Bajto\v{s} et al.}
{2020}
{Predictive blacklisting from 12M IDS alerts via sequential rule
mining.
\begin{itemize}[nosep]
  \item \acq{Channel~3 (SQ13).}
  \item \pro{Practical precedent for sequence mining on security data.}
  \item \con{Operates on alert streams, not log events:
  \textbf{outside thesis scope}.}
  \item \usmark{$\times$ Excluded: alerts are out of scope (EC1).}
\end{itemize}}

%---------------------------
\subsection{Log Representation (2 sources)}
%---------------------------

Wu et~al.~\cite{wu2023logrep} compare classical message-count vectors,
template-ID sequences, and frozen contextual embeddings as input
representations for downstream AD, showing that representation choice
interacts with dataset properties, directly informing D5.

\slrpaper{Wu, Li, Khomh (2023) -- On the Effectiveness of Log Representation for Log-Based AD}
{Xingfang Wu, Heng Li, Foutse Khomh}
{Empir.\ Softw.\ Eng.\ 28(6):137, 2023}
{Systematic comparison of 6 log representations across 7 ML models.
\begin{itemize}[nosep]
  \item \acq{Channel~3 (SQ10).}
  \item \datasets{HDFS, BGL, Thunderbird (in-corpus) and Spirit
  (out-of-corpus).}
  \item \pro{Shows representation choice interacts with dataset
  properties. No single representation is universally best: message-count vectors work
  best for sequential deep models, while BERT gives the strongest semantic
  representation.}
  \item \con{Does not link to dataset-centric D5 assessment.}
  \item \usmark{$\bigstar$ Primary: D5 representation sensitivity.}
\end{itemize}}

\slrpaper{Landauer, Wurzenberger et al.\ (2023) -- AMiner: Modular Log Data Analysis Pipeline}
{Max Landauer, Markus Wurzenberger et al.}
{2023}
{AECID-lineage modular pipeline.
\begin{itemize}[nosep]
  \item \acq{Channel~1 (seed).}
  \item \pro{Reference for reproducible modular pipeline.}
  \item \con{Tool-oriented.}
  \item \usmark{$\circ$ Contextual.}
\end{itemize}}

%---------------------------
\subsection{Datasets: Calibration, Primary, Legacy, Survey-Only (7 sources)}
%---------------------------

\slrpaper{Landauer, Skopik et al.\ (2021) -- Have It Your Way: AIT-LDS}
{Max Landauer, Florian Skopik, Markus Wurzenberger et al.}
{IEEE Trans.\ Reliab.\ 70(1):402--415, 2021}
{Model-driven simulation testbed for customisable log datasets.
\begin{itemize}[nosep]
  \item \acq{Channel~1 (seed).}
  \item \datasets{Introduces AIT-LDS (the Russell Mitchell, Santos and
  Wilson scenarios profiled and used for calibration in this thesis).}
  \item \pro{The calibration instrument: ground-truth properties are
  known and configurable.}
  \item \con{Simulated; transfer gap must be analyzed.}
  \item \usmark{$\bigstar$ Primary: calibration backbone.}
\end{itemize}}

\slrpaper{Landauer, Skopik et al.\ (2023) -- Maintainable Log Datasets for Evaluation of IDS (AIT-ADS)}
{Max Landauer, Florian Skopik, Maximilian Frank et al.}
{IEEE TDSC 20(4):3466--3482, 2023}
{Extends AIT-LDS toward reproducible maintenance and richer attack
coverage.
\begin{itemize}[nosep]
  \item \acq{Channel~1 (seed).}
  \item \datasets{Introduces AIT-ADS (maintainable extension of the
  AIT-LDS simulation lineage).}
  \item \pro{Second calibration dataset; cross-configuration
  robustness.}
  \item \con{Same simulation lineage.}
  \item \usmark{$\bigstar$ Primary: calibration.}
\end{itemize}}

\slrpaper{Zhu, He, He, Liu, Lyu (2023) -- Loghub}
{Jieming Zhu, Shilin He et al.}
{ISSRE 2023, pp.\ 355--366 (GL: github.com/logpai/loghub)}
{Aggregator of 19 log datasets (HDFS, Hadoop, OpenStack, BGL,
Thunderbird, \ldots).
\begin{itemize}[nosep]
  \item \acq{Channel~1 (foundational seed).}
  \item \datasets{Aggregates 19 log datasets, incl.\ HDFS, Hadoop,
  OpenStack, BGL, Thunderbird (the legacy-benchmark source for this
  thesis).}
  \item \pro{Reference for legacy-baseline datasets.}
  \item \con{Source of the benchmarks whose reuse the thesis
  critiques.}
  \item \usmark{$\bigstar$ Primary: legacy baseline source.}
\end{itemize}}

\slrpaper{Oliner, Stearley (2007) -- What Supercomputers Say}
{Adam Oliner, Jon Stearley}
{DSN 2007, pp.\ 575--584}
{Classic study of BGL, Red Storm, Thunderbird, Spirit, Liberty.
\begin{itemize}[nosep]
  \item \acq{Channel~2 (backward snowball from \cite{landauer2024}).}
  \item \datasets{BGL, Red Storm, Thunderbird, Spirit, Liberty (the
  original CFDR supercomputer logs; BGL and Thunderbird are in-corpus).}
  \item \pro{Provenance for BGL/Thunderbird (CFDR).}
  \item \con{Fault-diagnosis, not security: exactly the mismatch
  problem.}
  \item \usmark{$\bigstar$ Primary: legacy baseline provenance.}
\end{itemize}}

\slrpaper{Schroeder, Gibson (2006) -- A Large-Scale Study of Failures in HPC Systems}
{B.\ Schroeder, G.A.\ Gibson}
{DSN 2006, pp.\ 249--258}
{Failure characterization on HPC.
\begin{itemize}[nosep]
  \item \acq{Channel~2 (backward snowball from \cite{landauer2024}).}
  \item \datasets{LANL HPC failure records (22 systems, CFDR; reliability
  data, out-of-corpus).}
  \item \pro{Background on what BGL/TB labels actually represent
  (failures, not attacks).}
  \item \con{Not a security paper; underlines label-target mismatch.}
  \item \usmark{$\circ$ Contextual.}
\end{itemize}}

\slrpaper{Landauer, Hotwagner et al.\ (2026) -- CAM-LDS: Cyber Attack Manifestations}
{Max Landauer, Wolfgang Hotwagner et al.}
{2026}
{Dataset emphasising attack manifestations with semantic annotations.
\begin{itemize}[nosep]
  \item \acq{Channel~1 (seed).}
  \item \datasets{Introduces CAM-LDS (the seven attack-scenario groups
  S1--S7 profiled in this thesis).}
  \item \pro{Potential D5-rich extension dataset.}
  \item \con{Very recent; not yet widely validated.}
  \item \usmark{$\star$ Secondary: potential extension.}
\end{itemize}}

\slrpaper{NIST SP 800-92 (Kent, Souppaya, 2006) -- Guide to Computer Security Log Management}
{Karen Kent, Murugiah Souppaya}
{NIST SP 800-92, 2006 (GL)}
{Authoritative log-management terminology.
\begin{itemize}[nosep]
  \item \acq{Channel~3 (SQ13 / targeted search).}
  \item \pro{Anchor for log/telemetry/alerts definitions.}
  \item \con{Operational guide, not research.}
  \item \usmark{$\bigstar$ Primary: terminology anchor.}
\end{itemize}}

\textbf{Note:} Dataset-specific sources referenced in the thesis
(DARPA TC CADETS/THEIA, ATLASv2, UWF-ZeekData22, COMISET~\cite{comiset}, Mordor/OTRF) are
accessed through public repositories and are not counted in the
93-source total, as they entered the thesis through direct dataset
acquisition rather than the literature-review process.

%---------------------------
\subsection{SLR Methodology (3 sources)}
%---------------------------

\slrpaper{Kitchenham, Charters (2007) -- Guidelines for Performing Systematic Literature Reviews in SE}
{Barbara Kitchenham, Stuart Charters}
{EBSE Tech.\ Report EBSE-2007-01, 2007 (GL)}
{De-facto SLR guideline for SE.
\begin{itemize}[nosep]
  \item \acq{Channel~3 (SQ8).}
  \item \pro{Structures the peer-reviewed search track.}
  \item \con{Does not address grey literature.}
  \item \usmark{$\bigstar$ Primary: SLR methodology.}
\end{itemize}}

\slrpaper{Page et al.\ (2021) -- PRISMA 2020}
{Matthew J.\ Page, Joanne E.\ McKenzie et al.}
{BMJ 372:n71, 2021}
{Reporting standard for systematic reviews.
\begin{itemize}[nosep]
  \item \acq{Channel~3 (SQ8).}
  \item \pro{Reporting template for the flow diagram.}
  \item \con{Medical lineage; SE adaptation via Kitchenham needed.}
  \item \usmark{$\bigstar$ Primary: reporting template.}
\end{itemize}}

\slrpaper{Garousi, Felderer, M\"antyl\"a (2019) -- Guidelines for Including Grey Literature and Conducting MLRs in SE}
{Vahid Garousi, Michael Felderer, Mika V.\ M\"antyl\"a}
{Inf.\ Softw.\ Technol.\ 106:101--121, 2019}
{Principled grey-literature inclusion.
\begin{itemize}[nosep]
  \item \acq{Channel~3 (SQ8).}
  \item \pro{Enables formal inclusion of MITRE ATT\&CK, dataset
  READMEs, tool docs.}
  \item \con{Grey-literature criteria inherently subjective; requires
  labeling scheme.}
  \item \usmark{$\bigstar$ Primary: MLR methodology.}
\end{itemize}}



%------------------------------------------------------------------
\section{Update-round additions}
These four sources were admitted in the update round (the update-round section of the thesis); they are counted in the dataset-usage record but not in the systematic-search statistics.

\slrpaper{Yan, Shi, Ren et al.\ (2026) -- Log Anomaly Detection via Transformers Pre-Trained on Massive Unlabeled Data}
{Senming Yan, Lei Shi, Jing Ren, Wei Wang, Limin Sun, Wei Zhang}
{IEEE Trans.\ Netw.\ Sci.\ Eng.\ 13:5943--5960, 2026}
{Transformer detector pre-trained on large unlabeled log datasets, then
fine-tuned for anomaly detection.
\begin{itemize}[nosep]
  \item \acq{Update round (dataset-usage record).}
  \item \datasets{HDFS, BGL, Thunderbird (all in-corpus).}
  \item \usmark{$\circ$ Contextual: usage evidence for the census.}
\end{itemize}}

\slrpaper{Sedl\'a\v{c}ek, \v{Z}\'adn\'ik, Barto\v{s} (2025) -- Anomaly Detection in Log Data: A Comparative Study}
{Ond\v{r}ej Sedl\'a\v{c}ek, Martin \v{Z}\'adn\'ik, V\'aclav Barto\v{s}}
{Int.\ Conf.\ Netw.\ Serv.\ Manag.\ (CNSM) 2025}
{Comparative evaluation of log-AD methods on the canonical benchmarks.
\begin{itemize}[nosep]
  \item \acq{Update round (dataset-usage record).}
  \item \datasets{HDFS, BGL, Thunderbird (all in-corpus).}
  \item \usmark{$\circ$ Contextual: usage evidence for the census.}
\end{itemize}}

\slrpaper{Xu, Huang, Fox, Patterson, Jordan (2009) -- Detecting Large-Scale System Problems by Mining Console Logs}
{Wei Xu, Ling Huang, Armando Fox, David Patterson, Michael Jordan}
{SOSP 2009}
{The origin paper of the HDFS benchmark: console-log mining with
PCA-based detection on Hadoop File System logs (and Darkstar game-server
logs). Notable for the census: the paper that \emph{introduced} the
HDFS dataset also evaluates anomaly detection on it, so HDFS's career as
an AD benchmark begins with its own release.
\begin{itemize}[nosep]
  \item \acq{Update round (dataset-usage record).}
  \item \datasets{HDFS (in-corpus); Darkstar (out of corpus). The
  ``Hadoop File System logs'' are the HDFS dataset, not the Loghub
  YARN-container Hadoop dataset.}
  \item \usmark{$\circ$ Contextual: historically foundational usage
  evidence.}
\end{itemize}}

\slrpaper{Zhang, Chen, Que et al.\ (2025) -- LogPurge: Log Data Purification for Anomaly Detection via Rule-Enhanced Filtering}
{Shenglin Zhang, Ziang Chen, Zijing Que, Yilun Liu, Yongqian Sun, Sicheng Wei, Dan Pei, Hailin Li}
{arXiv:2511.14062, 2025}
{Rule-enhanced filtering to purify log data before anomaly detection.
\begin{itemize}[nosep]
  \item \acq{Update round (dataset-usage record).}
  \item \datasets{BGL (in-corpus), Liberty (out of corpus,
  the section on datasets not profiled), and an industrial dataset.}
  \item \usmark{$\circ$ Contextual: usage evidence for the census.}
\end{itemize}}
